\documentclass{ximera}

% vier voorkeuren die je meteen zelf kan instellen:
\def\uitbr{0} % waarde 1 als je Uitbreiding in de marge wilt, 0 als je dat niet wilt 
\def\wsg{0} % waarde 1 als je verwijzingen naar Wiskunde Samen gevat² in de marge wilt, 0 als je dat niet wilt
\def\rectoverso{0} % waarde 1 als je het PDF-bestand recto-verso wilt laten afdrukken, 0 als je het recto wilt
\def\voetn{1} % waarde 1 als je de voetnoten wilt, 0 als je dat niet wilt 

% marges:
\usepackage[paper=a4paper,margin=3.4cm,marginparwidth=2cm]{geometry}

% packages algemeen:
\pdfOnly{
\usepackage[english,dutch]{babel}
}
\usepackage{amsfonts}
\usepackage{amsmath} %  o.a. \eqref en \text, omgeving equation*
\usepackage{amssymb} % o.a. \nmid
\usepackage{amsthm} % o.a. omgeving proof
\usepackage{graphicx} % o.a. figuren
\usepackage{xcolor} % \color
\usepackage[disable]{todonotes} % \todo 
\usepackage{comment} % omgeving comment
\usepackage{xcolor} % kleuren 
\usepackage[skip=7pt plus1pt, indent=0pt]{parskip} % skip = verticale ruimte tussen twee alinea's, indent = insprong bij nieuwe alinea
\usepackage{multicol} % kolommen
\usepackage{xfrac} % breuken

\usepackage{siunitx} % SI eenheden
\usepackage{eso-pic} % achtergrond voorpagina
\usepackage{tabularx} % kolommen voor rekenmachine
\usepackage{rotating} % voor omgeving turn

% voor figuren met PSTricks:
\usepackage{pstricks} 
\usepackage{pstricks-add}
\usepackage{pst-plot}
\usepackage{pst-node}
\usepackage{pst-coil}
\usepackage{auto-pst-pdf}

% voor figuren met TikZ:
\usepackage{tikz} 
\usepackage{tkz-euclide} 
\usepackage{pgfplots} 
\usetikzlibrary{calc,intersections,through,backgrounds,patterns} 
\pgfplotsset{compat=newest}
\usepgfplotslibrary{fillbetween,colormaps}
\usepackage{import}
\usepackage{tikz-3dplot}

% zelfde parskip in minipage:
\makeatletter
\setlength{\parskip}{\medskipamount}
\newcommand{\@minipagerestore}{\setlength{\parskip}{\medskipamount}}
\makeatother

% voor veranderen van de naam Bibliografie naar Referentielijst:
\pdfOnly{
\addto\captionsdutch{\renewcommand{\bibname}{Referentielijst}}

% voor trefwoordenlijst:
\usepackage{imakeidx}
\makeindex[title=Trefwoordenlijst,program=makeindex,options=-s index.tex,columns=2,intoc=true]
}
% voor nummeren van vergelijkingen:
%WIM%\numberwithin{equation}{chapter} % als je dit desactiveert dan worden vergelijkingen in bijvoorbeeld hoofdstuk 3 genummerd als (1), (2) etc. in plaats van (3.1), (3.2) etc.

% voor hyperlinks en bladwijzers in PDF-bestand:
%WIM%\usepackage[bookmarksopen,bookmarksopenlevel=0,hypertexnames=false,pdfa,bookmarksnumbered]{hyperref} 
%WIM%\usepackage{bookmark} 

% voor hyperlinks: 
\hypersetup{pdfborder={0 0 0}, pdfstartpage=1,linkbordercolor={1 0 0}}%pdfborder={0 0 0} is geen rand rond hyperlinks, pdfborder={1 1 1} wel

% korter commando voor displaystyle, om bijvoorbeeld breuken groter te drukken in doorlopende tekst: 
\newcommand{\D}{\displaystyle}

% (veel)gebruikte verzamelingen:
\newcommand\NN{\mathbb{N}} % verzameling van de natuurlijke getallen
\newcommand\QQ{\mathbb{Q}} % verzameling van de rationale getallen
\newcommand\RR{\mathbb{R}} % verzameling van de re\"ele getallen
\newcommand\ZZ{\mathbb{Z}} % verzameling van de gehele getallen

% (veel)gebruikte operatoren:
\def\co{\operatorname{co}} % coördinaat van een punt
\def\ggd{\operatorname{ggd}} % positieve grootste gemene deler van twee gehele getallen niet beide nul
\def\gr{\operatorname{gr}} % graad van een veelterm

% voor kleur grafieken (donkergroen):
\definecolor{graf}{RGB}{0,100,0} 

% voor lijsten: 
% \usepackage{enumerate}
%WIM%\usepackage[shortlabels]{enumitem} 
\setlist{topsep=0em, itemsep=-0.15em}

% voor small bullet:
\newcommand\sbullet[1][.5]{\mathbin{\vcenter{\hbox{\scalebox{#1}{$\bullet$}}}}}

% voor meer verticale ruimte tussen vergelijkingen in align
\addtolength{\jot}{0.1cm}

% voor meervoudige voetnoten:
\usepackage{fnpct}

% voor het onderdrukken van voetnoten als optie:
\usepackage{letltxmacro}
\LetLtxMacro\Oldfootnote\footnote
\newcommand{\EnableFootNotes}{%
  \LetLtxMacro\footnote\Oldfootnote%
}
\newcommand{\DisableFootNotes}{%
  \renewcommand{\footnote}[2][]{\relax}
}

% voor verticale lijn in de marge (uitbreiding):
\usepackage[framemethod=default]{mdframed}
\usepackage{marginnote}
\reversemarginpar
\ifthenelse{\uitbr < 1}{\definecolor{rood}{RGB}{255,255,255}}{\definecolor{rood}{RGB}{254,64,64}} %HEX: #fe4040
\mdfdefinestyle{uitbreiding}{%
    topline=false,
    rightline=false,
    bottomline=false,
    leftline=true,
    linecolor=rood,
    linewidth=5pt,
    rightmargin=0pt,
    skipabove=10pt,% ipv 3
    skipbelow=0pt,
    leftmargin=-25pt,
    innerleftmargin=20pt,
    innerrightmargin=0pt,
    innertopmargin=0pt,
    innerbottommargin=0pt%,
%	needspace=30pt %minimumhoogte vooraleer lijn wordt gesplitst
    }
\newenvironment{Uitbreiding}{
    \marginpar{
        \center
		\vspace{0.1cm}
		\vspace{7pt}
		\rotatebox{90}{\color{rood}\Large \bf Uitbreiding}
	}
    \begin{mdframed}[style=uitbreiding]
    }{\vspace{-0.05cm}
    \end{mdframed}
}


% omgevingen voor lemma, definitie, voorbeeld etc.:
\newtheoremstyle{mystyle}
    {0em} % Space above
    {0em} % Space below
    {\itshape} % Body font
    {} % Indent amount
    {\bfseries} % Theorem head font
    {.} % Punctuation after theorem head
    {.5em} % Space after theorem head
    {} % Theorem head spec (can be left empty, meaning `normal')
\theoremstyle{mystyle}
%WIM%\newtheorem{lemma}{Lemma}[chapter] % als je [chapter] desactiveert dan Voorbeeld 3 in plaats van Voorbeeld 5.3, als je [chapter] vervangt door [section] dan Voorbeeld 5.2.3 in plaats van Voorbeeld 5.3 
%\theoremstyle{definition} % als je dit activeert dan is wat in de omgeving staat niet cursief gedrukt
\newtheorem{oefening}[lemma]{Oefening}
\newtheorem{definitie}[lemma]{Definitie} 
\newtheorem{voorbeeld}[lemma]{Voorbeeld}
\newtheorem{eigenschap}[lemma]{Eigenschap} 
\newtheorem{stelling}[lemma]{Stelling} 
\newtheorem{gevolg}[lemma]{Gevolg} 
\newtheorem{afspraak}[lemma]{Afspraak} 
\newtheorem{werkwijze}[lemma]{Werkwijze}

% omgeving proof, aangepaste ruimte: 
\makeatletter
\renewenvironment{proof}[1][\proofname]{\par
  \vspace{-\topsep}% remove the space after the theorem
  \pushQED{\qed}%
  \normalfont
  \topsep0pt \partopsep0pt % no space before
  \trivlist
  \item[\hskip\labelsep
        \itshape
    #1\@addpunct{.}]\ignorespaces
}{%
  \popQED\endtrivlist\@endpefalse
  \addvspace{0pt plus 0pt} % no space after
}
\makeatother

% kaderstijlen uit SOHO Wiskunde Plantyn:
\colorlet{steunkleur}{black}
\colorlet{steunkleurlicht}{steunkleur!30!white}
\colorlet{steunkleurkader}{steunkleur!7!white}
\colorlet{steunkleurkaderlicht}{steunkleur!2!white}

\mdfdefinestyle{kaderstijl_vol_licht}
{skipabove=6pt, 
skipbelow=6pt, 
backgroundcolor=steunkleurkader,
linecolor=steunkleurlicht,
linewidth = 0.4pt, 
topline=true,
bottomline=true, 
rightline=true,
innerleftmargin=5pt,
innerrightmargin=5pt,
innertopmargin=5pt,
leftmargin=0cm,
rightmargin=0cm,
innerbottommargin=5pt,
needspace=30pt % minimumhoogte voor splitsen kader
}
\surroundwithmdframed[style=kaderstijl_vol_licht]{definitie}
\surroundwithmdframed[style=kaderstijl_vol_licht]{stelling}
\surroundwithmdframed[style=kaderstijl_vol_licht]{lemma}
\surroundwithmdframed[style=kaderstijl_vol_licht]{eigenschap}
\surroundwithmdframed[style=kaderstijl_vol_licht]{gevolg}

% voor omgevingen voor oefening en antwoord:
\newenvironment{Oefening}{%
    \begin{enumerate}[ 
    series=Oef,
    resume=Oef,
    leftmargin=1.78em,
    label={\bfseries\arabic*.},
    ref=\arabic*
    ]
    \item %
    }{%
    \end{enumerate}
}
\newenvironment{Antwoord}{%
    \begin{enumerate}[%
    series=Antw,
    resume=Antw,
    leftmargin=1.78em,
    label={\bfseries\arabic*.},
    ref=\arabic*
    ]
    \item %
    }{%
    \end{enumerate}
}

% als de omgeving Uitbreiding start met een kaderomgeving (definitie, stelling, eigenschap, lemma of gevolg) dan moet wat extra verticale ruimte voorzien worden, met het commando \uitbreidingstartmetkader:
\def\uitbreidingstartmetkader{\mbox{}\vspace{-0.205cm}}

% % voor schema van de staartdeling:
% \usepackage{stackengine}
% \setstackgap{S}{5pt}
% \stackMath\def\stackalignment{r}
% \newcommand{\myRule}[3][white]{\textcolor{#1}{\rule{#2}{#3}}}
% \let\ph\phantom % enkel voor tekst
% \newcommand{\mph}[1]{% enkel voor math mode
%     \mathcolor{white}{#1}%
% }
% \def\staartmin{\rule{0.25cm}{0.1mm}\myRule{0.3cm}{0.1mm}}
% \def\staartphmin{\myRule{0.25cm}{0.1mm}\myRule{0.3cm}{0.1mm}}
% \newcommand{\staartstreep}[1]{\rule{\widthof{$#1$}}{0.1mm}}
% \newcommand{\staartphstreep}[1]{\myRule{\widthof{$#1$}}{0.1mm}}

% voor kolommen met schema van Horner:
\newcommand{\kolbreed}{1.0cm}
\newcolumntype{H}{>{\centering\arraybackslash$} p{\kolbreed} <{$}}

% voor nieuw commando utikzdashed: onderlijn (zoals underline) maar dan met stippellijn:
\tikzset{
    cheating dash/.code args={on #1 off #2 ends #3}{%
        \csname tikz@addoption\endcsname{%
            \pgfgetpath\currentpath%
            \pgfprocessround{\currentpath}{\currentpath}%
            \csname pgf@decorate@parsesoftpath\endcsname{\currentpath}{\currentpath}%
            \pgfmathparse{max(#1-#3,0)}\let\dashphase=\pgfmathresult%
            \pgfmathparse{\csname pgf@decorate@totalpathlength\endcsname-#1+2*\dashphase}\let\rest=\pgfmathresult%
            \pgfmathparse{#1+#2}\let\onoff=\pgfmathresult%
            \pgfmathparse{max(floor(\rest/\onoff), 1)}\let\nfullonoff=\pgfmathresult%
            \pgfmathparse{max((\rest-\onoff*\nfullonoff)/\nfullonoff+#2, #2)}\let\offexpand=\pgfmathresult%
            \pgfsetdash{{#1}{\offexpand}}{\dashphase pt}}%
    },
    cheating dash per segment/.style args={on #1 off #2 ends #3}{
        /utils/exec=\csname tikz@options\endcsname,%
        decoration={show path construction,
            lineto code={\draw [cheating dash=on #1 off #2 ends #3] (\tikzinputsegmentfirst) -- (\tikzinputsegmentlast);},
            curveto code={\draw [cheating dash=on #1 off #2 ends #3] (\tikzinputsegmentfirst) .. controls (\tikzinputsegmentsupporta) and (\tikzinputsegmentsupportb) .. (\tikzinputsegmentlast);},
            closepath code={\draw [cheating dash=on #1 off #2 ends #3] (\tikzinputsegmentfirst) -- (\tikzinputsegmentlast);}
        },
        decorate,
    },
}
\newcommand{\utikzdash}[1]{%
    \tikz[baseline=(todotted.base)]{
        \node[inner sep=0pt,outer sep=1.5pt] (todotted) {#1};
        \draw[cheating dash per segment=on 2pt off 2pt ends 2pt, line width=0.4pt] (todotted.south west) -- (todotted.south east);
    }%
}%

\providecommand{\xmemph}[1]{\textit{#1}}

% voor nieuw commando underdashed: onderlijn met stippellijnen en ook het woord cursief zetten:
\newcommand{\underdashed}[1]{%
    {\em\utikzdash{\!#1\!}}%
}

% voor icoon TI-84 Plus met verwijzing naar filmpje:
\newcommand{\grmlink}{\raisebox{0cm}{\includegraphics[width=1cm]{TI84Plus-icoon}}}
\newcommand{\grmref}[1]{%
	\reversemarginpar%
	\marginpar{
		\vspace{-0.4cm}%
		\htmladdnormallink{\grmlink}{#1}
	}	
}

% GRM knoppen:
\newcommand{\GRM}[1]{\fbox{\rule[0mm]{0cm}{0.215cm}\textup{\texttt{#1}}}}
\newcommand{\wedgetext}{{\raisebox{0.02cm}{\begin{turn}{90}>\end{turn}}}}
\newcommand{\veetext}{\raisebox{0.2cm}{\begin{turn}{-90}>\end{turn}}}

% voor kolommen met GRM screens:
\newlength{\widthallscreens}
\newlength{\widthscreens}
\newlength{\spaceleftscreen}
\newlength{\spacerightscreen}
\newlength{\spacebetweenscreens}

\newcommand{\setscreens}{
	\setlength{\spaceleftscreen}{2pt}
	\setlength{\spacerightscreen}{2pt}
	\setlength{\spacebetweenscreens}{8pt}
	\addtolength{\linewidth}{-28pt} % 3*8pt tussen vier screens en 2 pt links en 2 pt rechts
	\setlength{\widthallscreens}{\linewidth}
	\setlength{\widthscreens}{0.25\widthallscreens}
	\addtolength{\linewidth}{28pt}
	\newcolumntype{G}{p{\widthscreens}}
	\newcolumntype{s}{p{\spacebetweenscreens}}
	\newcolumntype{L}{p{\spaceleftscreen}}
	\newcolumntype{R}{p{\spacerightscreen}}
	\setlength{\tabcolsep}{0pt}
}

% voor icoon met verwijzing naar Wiskunde Samen gevat² in de marge:
% \newcommand{\wsglink}{\raisebox{0cm}{\includegraphics[width=1cm]{wsglogo}}}
\newcommand{\wsglink}{\raisebox{0cm}{LOGO}}
\newcommand{\htmladdnormallink}[2]{\href{#2}{#1}}
\newcounter{pagnrwsg}
\newcommand{\wsgref}[3]{% 
    % #1 woord dat onderstippeld wordt
    % #2 pagina van Wiskunde Samen gevat² waar je op terecht komt als je op de link klikt  
    % #3 afgedrukt op het logo van Wiskunde Samen gevat² in de marge: één of meerdere paginanummers
    \ifthenelse{\wsg < 1}{#1}{%
        \underdashed{#1}%
        {\setcounter{pagnrwsg}{#2}}%
        {\addtocounter{pagnrwsg}{14}}%
        \reversemarginpar%
        \marginpar{\vspace{-0.6cm}%
           \htmladdnormallink{\wsglink}{https://online.fliphtml5.com/sanky/laea/\#p=\arabic{pagnrwsg}}%
            \raisebox{0.41cm}[0cm][0cm]{%
                \hspace{-1.5cm}\makebox[2cm][c]{\colorbox{white}{\texttt{\footnotesize{#3}}}}%
            }%
        }%
    }%
}

% voor invoegen blanco pagina bij de optie recto-verso:
\def\blancobijrectoverso{
	\ifthenelse{\rectoverso < 1}{\clearpage}{
	\clearpage
	\thispagestyle{empty}
	\mbox{}
	\clearpage
	}
}

% aantal pagina's van het bestand:
\ifthenelse{\rectoverso < 1}{\def\totpag{57}}{\def\totpag{66}}

% documenteigenschappen:
% \usepackage{hyperxmp}
% \hypersetup{
% pdftitle={Open Source Wiskunde Aan zet: Veeltermen}, 
% pdfnumpages={\totpag},
% pdfauthor={Koen De Naeghel},
% pdflang={nl},
% pdfkeywords={wiskunde, open source, wiskunde aan zet, veeltermen, secundair onderwijs, tweede graad, doorstroomfinaliteit},
% pdfsubject={Veeltermen},
% pdfcopyright={\unichar{"24B8} 2024 Koen De Naeghel},
% pdfdate={17 december 2024},
% pdfapart=1,
% pdfstartview=}

% voor het benoemen van titel, auteur en datum:
% \title{Veeltermen}
% \author{Auteur: Koen De Naeghel}
% \date{\today}

\newcommand\BackgroundPic{
    \put(-260,-125){
    \parbox[b][\paperheight]{\paperwidth}{%
    \vfill
    \centering
    \includegraphics[height=\paperwidth, keepaspectratio]{WaZlogo}%
    \vfill
}}}

\addPrintStyle{..}
\begin{document}
	\author{Wim Obbels}
	\xmtitle{De norm van een complex getal}{}
	\label{xim:complexe_getallen_norm}


\begin{xmuitweiding}[Over de absolute waarde van een reëel getal]
De absolute waarde $|a|$ van een reëel getal $a$ lijkt op het eerste zicht een tamelijk eenvoudig begrip, want $|a|$ is 'het getal $r$ zonder het teken'.
Op de getallenrechte kan je $|a|$ ook beschouwen als de afstand van $a$ tot de oorsprong, en $|a-b|$ als de afstand tussen $a$ en $b$.
In formules en berekeningen doet de absolute waarde soms vervelend, omdat de absolute waarde $|a+b|$ van een som niet kan geschreven worden als de som van de absolute waarden van de termen, dus als $|a|+|b|$.
\end{xmuitweiding}

%Voor complexe getallen is 'het getal zonder het teken' weinig zinvol, maar de interpretatie als afstand blijkt wel erg nuttig.
%In plaats van over 'absolute waarde' spreekt men bij complexe getallen over de modulus of de norm, en dat zijn eigenlijk andere namen voor 'afstand' bij complexe getallen.


De \textit{absolute waarde} $|a|$ van een reëel getal $a\in\R$ geeft de afstand weer van dit punt tot de oorsprong op de reële rechte. Voor complexe getallen neemt men de afstand tot de oorsprong van het complexe vlak.

% Om het begrip \textit{absolute waarde} $|a|$ van een reëel getal $a\in\R$ te veralgemenen naar complexe getallen, begrijp je $|a|$ best als de afstand van het punt $a$ op de reële rechte tot de oorsprong. En bij complexe getallen gebruiken we niet het woord 'absolute waarde', maar wel 'norm' of 'modulus': 


\begin{definition}[Modulus van een complex getal]\label{def:complex_toegevoegde} \nl 
	
De \textbf{modulus} (of \textbf{norm}) van een complex getal $z=a+bi$, genoteerd $|z|$, is het \textit{positief reëel} getal 
%
	\formulevb{|z| \perdef  |a+bi| \perdef \sqrt{a^2+b^2}}{|3+4i| = \sqrt{3^2+4^2} = \sqrt{9+16}=\sqrt{25}=5.}
%

\begin{image}[0.6\textwidth]
	\begin{tikzpicture}[scale=5]%,cap=round,transform canvas={scale=0.5}]
	
	\tikzmath{\hoek = 30;}
	
	\draw[->] (-0.1,0) -- (1.2,0) node[below left] {Re$(z)$};
	\draw[->] (0,-0.1) -- (0,0.7) node[below left] {Im$(z)$};
	
	\draw[color=blue,thick] (0:0)     node[below,black]{$0$} 
	                        --        node[above, rotate=\hoek] {\small $|z|=\sqrt{a^2+b^2}$} 
							(\hoek:1) node [yshift=1pt,above] {$z=a+bi$} ; 
    
	\draw[color=black] (\hoek:1) node[name=P,circle, fill=black, radius=1pt,scale=0.8] {} ;  

    \draw[dashed] ({cos(\hoek)},0) node[below]{$a$} -- (P);
	\draw[dashed] (0,{sin(\hoek)}) node[circle, fill=black, radius=1pt,scale=0.5] {} node[left] {$b$} -- (P);
	\draw[dashed] ({cos(\hoek)},0) node[circle, fill=black, radius=1pt,scale=0.5] {};

	\begin{scope}[xshift=1.5cm]
	
	\draw[color=blue,thick] 
					(0:0)      node[below,black]{$0$}
					--         node[above, rotate=\hoek] {\small $|z|=\sqrt{a^2+b^2}$} 
					(\hoek:1)  node[above,black]{$a+bi$}; 
    
	\draw[color=red] (0,0) -- node[below]{$a$} ({cos(\hoek)},0);  
	\draw[color=red] (\hoek:1) -- node[right]{$b$} ({cos(\hoek)},0);  


	\end{scope}
	
	\end{tikzpicture}
\end{image}

De modulus $|z|$ is meetkundig \textbf{de afstand van $z$ tot de oorsprong} (wegens Pythagoras). 
%en dus ook de \textit{lengte van de vector} $\vec{z}$.
\end{definition}
 
 
\begin{quickquestion*}{}

	Teken alle complexe getallen waarvoor \(|z| = 1\) in het complexe vlak. 
\end{quickquestion*}

% \begin{quickquestion}{}{}    % met een nummer ...
% 	Teken alle complexe getallen waarvoor \(|z| = 1\) in het complexe vlak. 
% \end{quickquestion}

\begin{example}\nl
	\begin{xmmulticols}[4]
		\begin{question} $|i|      = \answer[onlineshowanswerbutton]{1}$        \end{question}
		\begin{question} $|-i|     = \answer[onlineshowanswerbutton]{1}$        \end{question}
		\begin{question} $|-2|     = \answer[onlineshowanswerbutton]{2}$        \end{question}
		\begin{question} $|2i|     = \answer[onlineshowanswerbutton]{2}$        \end{question}
		\begin{question} $|1-1|    = \answer[onlineshowanswerbutton]{0}$        \end{question}
		\begin{question} $|1|+|-1| = \answer[onlineshowanswerbutton]{2}$        \end{question}
		\begin{question} $|3-4i|   = \answer[onlineshowanswerbutton]{5}$        \end{question}
		\begin{question} $|1+i|    = \answer[onlineshowanswerbutton]{\sqrt{2}}$ \end{question}
	\end{xmmulticols}	
\end{example}


\todo{Video verhuizen naar in extra activity 'norm+modulus'}
De meetkundige betekenis van de modulus en complex toegevoegde wordt uitgelegd in onderstaande video:
	
 	\youtube{uSl9oLBjw1U}


De modulus heeft volgende basiseigenschappen, die je allemaal tamelijk eenvoudig kan bewijzen.

\begin{proposition}[Eigenschappen modulus]\label{eig:complexe_modulus}
	Voor complexe getallen $z,z_1,z_2\in \C$ geldt
	%\begin{xmmulticols}
	\begin{enumerate}
		\item $|z|$ is de afstand van $z$ tot de oorsprong.
		\item $|z_1-z_2|$ is de afstand tussen $z_1$ en $z_2$.
		\item $|z|= |-z|$.		
		\item $|z|=0 \iff  z=0 $.
		\item $|z_1\cdot z_2| = |z_1| \cdot |z_2|$.
%		\item $z \cdot \overline{z} = |z|^2$
%		\item $\dfrac{z \cdot \overline{z}}{|z|^2} = 1$ \label{eig:modulus:invers}
		
		%\item $\displaystyle \left| \frac{1}{z}\right|= \frac{1}{|z|}$
		%\item $\left| \displaystyle \frac{z_1}{z_2}\right|=\displaystyle \frac{|z_1|}{|z_2|}$
		%\item \important{|z_1+z_2| \leq |z_1|+|z_2|}
		
	\end{enumerate}
	%
	\todo{Driehoeksongelijkheid terug toevoegen}
	%\end{xmmulticols}
	LET OP: in het algemeen geldt niet $\xcancel{|z_1+z_2| = |z_1|+|z_2|}$ (want bv. $0= |1+(-1)| \neq |1|+|-1| = 2)$
	
\end{proposition}

% Om op één pagina te passen, is het handig de opgave's in de PDF niet te herhalen.
% TODO: misschien toch optioneel de oplossingen in de 'standaard' versie zetten!!!
% \pdfOnly{
% De bewijzen van deze eigenschappen zijn oefeningen. 
% }

\begin{onlineOnly}   % DOES NOT WORK ???
    
\begin{exercise} Bewijs alle uitspraken van Eigenschap~\href{eig:complexe_modulus}.

	\begin{question} 
		$|z|$ is de afstand tussen $z$ en de oorsprong in het complexe vlak.
		\begin{oplossing}
			Het bewijs staat al in Definitie~\ref{def:complex_toegevoegde}.
		\end{oplossing}
	\end{question}
	\begin{question} 
		$|z_1 - z_2|$ is de afstand tussen $z_1$ en $z_2$ in het complexe vlak.
		   
		\begin{oplossing}
			Wegens de parallellogramregel is  $z_1-z_2$  als vector gelijk aan de vector tussen $z_1$ en $z_2$, en dus is $|z_1-z_2|$ ook de afstand tussen $z_1$ en $z_2$.
			\begin{image}[0.4\textwidth]
			    \begin{tikzpicture}[scale=5]%,cap=round,transform canvas={scale=0.5}]
				   
				   \draw[->] (-0.6,0) -- (1.4,0) node[above] {Re$(z)$};
				   \draw[->] (0,-0.6) -- (0,0.6)   node[below right] {Im$(z)$};
				   
				   \draw[color=blue,thick] (0:0)  -- (1,0.3); 
				   \draw[color=blue,thick] (0:0)  -- (0.5,0.5); 
				   \draw[color=blue,dashed] (0:0)  -- (-0.5,-0.5); 			   
				   %
				   \draw[color=black] (1,0.3) node[name=Z1,circle, fill=black, radius=1pt,scale=0.8] {} node [yshift=2pt,above] {$z_1$} ;  
				   \draw[color=black] (0.5,0.5) node[name=Z2,circle, fill=black, radius=1pt,scale=0.8] {} node [yshift=2pt,above] {$z_2$} ;  
				   %
				   \draw[color=black] (-0.5,-0.5) node[name=mZ2,circle, fill=black, radius=1pt,scale=0.8] {} node [yshift=2pt,above left ] {$-z_2$} ;  
				   %
				   \draw[color=black] (0.5,-0.2) node[name=Zv,circle, fill=black, radius=1pt,scale=0.8] {} node [yshift=-2pt,below right] {$z_1-z_2$} ;  
				   %
				   \draw[red] (Z1) -- node[above right] {\small $|z_1-z_2|$} (Z2);
				   \draw[red] (0,0) -- node[right] {\small $|z_1-z_2|$} (Zv);
				   %
				   \draw[dashed] (mZ2) -- (Zv);
				   \draw[dashed] (Z1) -- (Zv);
				   
				\end{tikzpicture}   
			\end{image}
		\end{oplossing}
	\end{question}		   
	
	\begin{question} $|z|= |-z|$.
		\begin{oplossing} Stel $z=a+bi$, bereken beide uitdrukkingen, en stel vast dat ze aan elkaar gelijk zijn.
			\begin{itemize}
				\item $| z| = \sqrt{a^2+b^2}$
				\item $|-z| = \sqrt{(-a)^2+(-b)^2} = \sqrt{a^2+b^2}$ \qed
			\end{itemize}
		\end{oplossing}
	\end{question}
		
	\begin{question} $|z| = 0 \iff z = 0$.
		\begin{oplossing}
			Stel $z=a+bi$, en bewijs de equivalentie door twee implicaties aan te tonen:
			
			1. $\boxed{\Rightarrow}$ Als $|z| = 0$, dan is $a^2+b^2 = 0$,. Omdat kwadraten positief zijn, kan $a^2+b^2=0$  alleen als $a=b=0$, en dus $z = 0$.
			
			2. $\boxed{\Leftarrow}$ Als $z = 0$, dan is $a=b=0$, en dus $a^2+b^2 = 0$, en dus $|z| = 0$.
			
		\end{oplossing}
	\end{question}
	
	\begin{question}  Bewijs dat $|z_1\cdot z_2| = |z_1| \cdot |z_2|$.
		\begin{hint} Bereken het linkerlid en het rechterlid en stel vast dat beide aan elkaar gelijk zijn.
		\end{hint}
		\begin{oplossing}
			Stel $z_1=a_1+b_1i$ en $z_2=a_2+b_2i$. Dan geldt
			% Mmm, \begin{align*} does not work with onlineOnly ... ???
			$$
			\begin{array}{rl}
			|z_1\cdot z_2| & = | (a_1+b_1i)(a_2+b_2i) | = |  (a_1a_2 - b_1b_2) + (a_1b_2+a_2b_1)i| \\
			& = \sqrt{ (a_1a_2 - b_1b_2)^2 + (a_1b_2+a_2b_1)^2}  \\
			&= \sqrt{ a_1^2a_2^2  + b_1^2b_2^2 -2a_1a_2b_1b_2 + a_1^2b_2^2 + a_2^2b_1^2 + 2a_1a_2b_1b_2}  \\
			&= \sqrt{ a_1^2a_2^2  + b_1^2b_2^2 + a_1^2b_2^2 + a_2^2b_1^2 }  \\
			\\
			|z_1|\cdot| z_2| & = \sqrt{a_2^2+b_1^2}\sqrt{a_2^2+b_2^2}= \sqrt{(a_2^2+b_1^2)(a_2^2+b_2^2)} \\
			&= \sqrt{ a_1^2a_2^2  + b_1^2b_2^2 + a_1^2b_2^2 + a_2^2b_1^2 }  
			\end{array}
			$$
			Merk op: je zou de notatie wat kunnen vereenvoudigen door te bewijzen dat $|z_1\cdot z_2|^2 = |z_1|^2 \cdot |z_2|^2$. Omdat de moduli toch altijd positief zijn, is dit voldoende.
		\end{oplossing}
	\end{question}
\end{exercise}


\end{onlineOnly}


\end{document}

